\documentclass[11pt,a4paper]{article}
\usepackage[margin=0.8in]{geometry}
\usepackage{hyperref}
\usepackage{enumitem}
\usepackage{titlesec}
\usepackage{tabularx}
\usepackage{array}
\usepackage{booktabs}

% Better section formatting
\titleformat{\section}{\large\bfseries}{}{0em}{}[\titlerule]

% Fix spacing after section titles
\titlespacing{\section}{0pt}{12pt}{6pt}

% Fix hyperref formatting
\hypersetup{
  colorlinks=true,
  linkcolor=blue,
  urlcolor=blue,
}

\begin{document}

% Header
\begin{center}
{\LARGE \textbf{FEMMInterpreter}} \\[4pt]
{\large Top-Level Configuration Attributes}
\end{center}

\section*{Magnetic Attributes}

\begin{table}[h]
  \centering
  \renewcommand{\arraystretch}{1.5}
  \begin{tabular}{>{\raggedright\arraybackslash}p{4.5cm} 
                  >{\raggedright\arraybackslash}p{3cm}
                  >{\raggedright\arraybackslash}p{9.5cm}}
    \toprule
    \textbf{Attribute} & \textbf{Type} & \textbf{Description} \\
    \midrule
    format\_version & float & The FEMM file format version number. \\
    frequency\_hz & float & Operating frequency in Hertz for AC analysis. \\
    solver\_precision & float & Numerical precision tolerance for the solver. \\
    min\_angle\_deg & float & Minimum angle in degrees for meshing constraints. \\
    model\_depth & float & Depth of the model in the z-direction (for 2D axisymmetric). \\
    length\_unit & str & Unit system for length measurements (e.g., `millimeters', `inches'). \\
    problem\_type & str & Type of analysis (e.g., `planar', `axisymmetric'). \\
    coordinate\_system & str & Coordinate system used (e.g., `cartesian', `polar'). \\
    comment\_text & str & User-defined metadata or notes about the model. \\
    \bottomrule
  \end{tabular}
\end{table}

\end{document}